\section{Introduction}\label{sec:intro}

\begin{quote}
"System 1 operates automatically and quickly, with little or no effort and no sense of voluntary control. System 2 allocates attention to the effortful mental activities that demand it, including complex computations."
\hfill\null 
\par\nopagebreak
\hfill --- Daniel Kahneman, \textit{Thinking, Fast and Slow}
\end{quote}



The field of video avatars~\cite{he2023gaia,tian2024emo,xu2024hallo,wang2024vexpress,chen2024echomimic,xu2024vasa,stypulkowski2024diffused,jiang2025loopy,lin2025cyberhost,gan2025omniavatar,kong2025multitalk,wang2025fantasytalking,hu2024animate,lin2025omnihuman,qiu2025skyreels} aims to construct models capable of synthesizing realistic videos of characters from human-centric driving signals. The ultimate goal can be seen as creating lifelike avatars that are virtually indistinguishable from humans, capable of demonstrating both reasoned action and authentic emotion. This domain has witnessed rapid advancements over the past few years, with model capabilities evolving from early lip movement synthesis~\cite{jiang2024mobileportrait,zhang2023sadtalker,facev2v,zhao2022tps,fomm,mraa} and portrait animation~\cite{jiang2025loopy,tian2024emo,fada,xu2024vasa,xu2024hallo,hallo3} to more recent half-body~\cite{lin2025cyberhost,emo2} and full-body generation~\cite{lin2025omnihuman}. As the scope of controllable generation expands, these models are increasingly expected to simulate a broader spectrum of human behaviors, moving beyond mere physical likeness to capture a character's authentic essence, as illustrated in Figure~\ref{fig:teaser}.


Recently, a series of audio-driven methods~\cite{lin2025omnihuman,kong2025multitalk,gan2025omniavatar,wang2025fantasytalking, qiu2025skyreels} based on Diffusion Transformers (DiT)~\cite{dit,esser2024sd3,seawead2025seaweed,kong2024hunyuanvideo,wan2025} have emerged, enabling the generation of human motion videos that synchronize with audio in specific contexts. These approaches typically follow a similar paradigm: they first pre-train a foundational video generation model~\cite{wan2025,seawead2025seaweed} and then introduce conditional audio inputs to achieve animation. While this yields reasonable results, a closer inspection reveals that these models often only capture the direct and simplistic correlations between the audio signal and the resulting motion. Consequently, they tend to generate only synchronized lip movements and simple, repetitive accompanying gestures. When judged on their potential to function as convincing avatars, these outputs still exhibit a significant gap in naturalness and plausibility when compared to authentic human behavior.


To understand the source of this gap, it is instructive to consider a prominent theory in human cognition, which posits that behavior is governed by two distinct systems~\cite{kahneman2011thinking,kahneman2013prospect}, often termed System 1 and System 2. System 1 is described as fast, unconscious, and reactive. This operational mode is analogous to that of current video avatar models, which excel at mapping input signals like audio directly to corresponding lip movements and simple gestures. In contrast, System 2 is characterized as deliberative, analytical, and effortful, involving reasoning about goals and context to deduce a course of action. This latter capability, the hallmark of intelligent behavior, remains challenging for existing methods to emulate robustly. Inspired by this dual-process model, we identify a clear path forward: leveraging the powerful reasoning capabilities~\cite{wei2022chain, wangself, yao2023tree, schick2023toolformer, park2023generative, yuan2024mora, wu2025automated} of modern Multimodal Large Language Models (MLLMs)~\cite{team2023gemini,GPT4o} to explicitly simulate the deliberative processes of System 2. We therefore propose a novel framework that integrates the principles of both systems, using MLLMs to simulate the goal-oriented aspects of System 2 while preserving the reactive capabilities analogous to System 1.


However, integrating the textual guidance from an MLLM into an avatar generation framework is non-trivial. Leveraging Chain-of-Thought (CoT)~\cite{wei2022chain} prompting, MLLMs articulate their reasoning process as explicit text, a natural medium for expressing logic. Consequently, introducing an MLLM to enhance the high-level semantic coherence of motion inevitably strengthens the role of text as a conditional input. This creates a potential conflict of modalities within existing avatar models, where the audio signal already governs lip-sync and rhythmic motion, and the reference image constrains the character's appearance and potential range of motion. These distinct inputs are all intricately linked to the final generated motion. Therefore, we argue that a novel framework must be designed to mitigate modal conflicts while effectively leveraging the interdependencies among these diverse conditions. Developing such a framework is the key to simultaneously simulating both ``System 1'' and ``System 2''.



Based on the preceding analysis, we propose a video avatar framework built upon a Multimodal Diffusion Transformer that features two key designs. First, to ensure a comprehensive simulation of both System 1 and System 2, we introduce an agentic system powered by MLLMs. These agents reason over multimodal inputs (text, reference image, audio) to generate high-level semantic context, providing a long-term, logically coherent signal (for ``System 2'') that complements the short-term, reactive audio signals (for ``System 1'').
Second, to effectively fuse these input conditions while mitigating modality interference, we employ dedicated Multimodal Branches for audio and text feature extraction, along with a Multimodal Attention mechanism for joint modeling of audio, text and video. This design iteratively updates and aligns audio and text features into a common semantic space. Furthermore, we introduce a distinct identity preservation strategy that avoids conditioning on the reference image during training. Instead, we guide the model by probabilistically conditioning on start/end frames during training and treating the reference image as a pseudo last frame at inference, which prevents the static image from interfering with dynamic, content-driven motion.
As a result of these designs, our model generates vivid and lifelike human motion from audio and a single image. The resulting videos exhibit remarkable contextual and semantic coherence, effectively simulating both the reactive and deliberative aspects of human behavior.

We summarize our main contributions as follows:
\begin{itemize}
    \item \textbf{A New Perspective on Avatar Modeling:} We introduce a new perspective for analyzing video avatars, framing the problem through the cognitive science lens of System 1 and System 2 thinking. Observing that current models primarily simulate System 1, we are the first to propose a holistic approach that models both.
    \item \textbf{A Framework for Dual-System Simulation:} To implement this vision, we propose a novel framework featuring two core components. First, MLLM-based agents generate deliberative guidance (``System 2''). Second, a specialized MMDiT architecture, equipped with a symmetric audio branch and our pseudo-last-frame strategy, synergistically fuses this guidance with reactive signals (``System 1'') , thereby resolving critical modal conflicts.
    \item \textbf{Strong Empirical Performance and Generalization:} Our method not only achieves highly competitive results on standard benchmarks but is also significantly preferred in user studies for its contextual naturalness and plausibility. Its versatility is further demonstrated by its successful extension to complex multi-person and non-human scenarios.
\end{itemize}


